\section{Introduction}
\label{sec:introduction}

Coastal lagoons are among the most productive coastal environments on Earth.
Their shallow, semi-enclosed geometry concentrates fisheries and biodiversity
along roughly 13\% of the global shoreline \citep{Carrasco2016, Kjerfve1994},
while making these systems acutely sensitive to the atmospheric and tidal
forcing that drives water exchange and renewal. In the Mediterranean basin,
where tidal amplitudes rarely exceed 0.3~m, the balance among tidal
oscillations, wind stress, and seiche phenomena varies considerably among
systems, yet wind consistently emerges as a significant driver of horizontal
exchange and a key modulator of the residence times and salinity gradients
that condition benthic habitat \citep{Niedda2007, Umgiesser2014}. The presence
of dense submerged vegetation, principally \emph{Posidonia oceanica} (L.)
Delile and \emph{Cymodocea nodosa} (Ucria) Ascherson along Mediterranean
shores, introduces a further layer of complexity. Seagrass canopies attenuate
incident wave energy \citep{Infantes2012} and extract momentum from the mean
flow through canopy drag, modifying near-bed turbulence and the vertical
velocity structure \citep{Nepf2012}. As global seagrass extent continues to
decline at accelerating rates \citep{Waycott2009}, understanding how
circulation controls residence times, dispersal, and habitat connectivity in
vegetated coastal systems has become an urgent scientific and management
priority. Addressing it requires circulation models that accurately reproduce
transport pathways rather than water levels alone, and this distinction makes
Lagrangian validation essential. By tracking the cumulative separation between
simulated and observed drifter trajectories over time, the Lagrangian skill
score provides a direct and discriminating measure of model accuracy that
conventional scalar metrics cannot replicate
\citep{Liu2011, Revelard2021}.

The \emph{Stagnone di Marsala}, on the north-western coast of Sicily, offers a
natural laboratory for these questions. The lagoon extends approximately 11~km
in the north--south direction and 2.5~km east--west, enclosing a surface area
of roughly 2200~ha between the mainland and the elongated Isola Grande
\citep{DeMarchis2012}. Its mean water depth is approximately 0.95~m, ranging
from a few tens of centimetres over the northern intertidal flats to a maximum
of about 3~m near the southern inlet, placing it firmly in the class of very
shallow coastal systems in which the depth-averaging assumption is, a priori,
difficult to justify. Hydrodynamic exchange with the open Mediterranean is
mediated by two mouths of strongly contrasting geometry: a narrow northern
inlet ($\sim$400~m wide, 0.30--0.40~m deep) whose shallow sill severely
restricts tidal exchange, partially mitigated by a historically dredged
channel, and a wide southern inlet ($\sim$2900~m wide, 1.0--1.5~m deep) that
controls the bulk of mass transfer
\citep{DeMarchis2012, CiraoloDeMarchis2009}. The lagoon supports no
significant freshwater input, and its bottom is colonised by
\emph{P.\ oceanica} in the central and southern basin and by
\emph{C.\ nodosa} in the shallower northern reaches, with shoot densities
reaching 1000--1200 plants~\si{\per\metre\squared} \citep{Ciraolo2009a}.
Designated as a \emph{Riserva Naturale Orientata} and a Natura 2000 site, the
Stagnone hosts recognised populations of protected flora and fauna. Despite
this status, satellite monitoring over 2003--2024 has documented a 75\%
reduction in the extent of the \emph{P.\ oceanica} atoll and cordon formations
\citep{Maltese2025}, underscoring the need to understand the hydrodynamic
processes that govern conditions within this enclosed basin.

The tidal regime is microtidal, with free-surface amplitudes of the order of
0.25~m dominated by semi-diurnal harmonics \citep{DiMarca2009}. In the absence
of riverine input and with weak tidal forcing, circulation is driven almost
exclusively by tidal oscillations and wind stress. Field experiments and 3D
numerical modelling conducted by \citet{DeMarchis2012} during a July 2006
campaign showed that tidal action dominates north--south exchange, while wind
is the controlling agent of east--west currents and, critically, generates
pronounced vertical recirculation structures entirely absent in tide-only
simulations. This established that vertical velocity shear between wind-driven
surface drift and near-bed return flow is quantitatively significant even at
sub-metre depths, and that the tendency towards hypersaline conditions in
summer \citep{LaLoggia2004} adds a baroclinic component that reinforces
stratification in the vertical.

Prior numerical investigations of the Stagnone have employed both
depth-averaged and three-dimensional approaches, each appropriate to the
research objective pursued. \citet{LaLoggia2004} and \citet{DiMarca2009}
applied 2D depth-averaged models to characterise tidal flushing, residence
times, and turbidity dynamics, yielding useful insight into the basin-scale
water budget and the role of biological productivity in light attenuation.
\citet{Ciraolo2009a} demonstrated through particle tracking the sensitivity of
simulated trajectories to the assumed velocity field, highlighting the
importance of current accuracy for transport prediction.
\citet{DeMarchis2012} showed that resolving the vertical dimension is
necessary to reproduce wind-driven recirculation. Their 3D non-hydrostatic
finite-volume model, validated against simultaneous velocity and water-level
measurements at multiple stations, achieved substantially better agreement
with observed east--west currents than any depth-averaged simulation. More
recently, \citet{Ingrassia2024} returned to a 2D unstructured-mesh framework
(MIKE 21) and conducted systematic sensitivity tests of bed friction, varying
the Gauckler--Strickler coefficient uniformly across the lagoon for different
vegetation scenarios. The value $K_s = 20$~m$^{1/3}$~s$^{-1}$ was identified
as most representative for the vegetated basin, yielding Nash--Sutcliffe
efficiencies up to 0.92 for water level. The authors noted that the 2D
approximation remains adequate for large-scale water-level budgeting while
acknowledging reduced predictive skill for horizontal velocity, particularly
in the northern, shallower sub-basin.

Building on these foundations, the present study addresses two aspects that
have not yet been incorporated jointly into a Stagnone circulation model. The
first is wave--current coupling. Spectral ocean waves generated in the western
Sicilian channel can penetrate both inlets and interact with the mean flow,
altering near-bed turbulence and the effective drag experienced by the
seagrass canopy. Coupled wave and flow modelling in comparably shallow
Mediterranean basins has shown this to be non-negligible. In the Venice Lagoon
\citet{Carniello2011} found the wave-induced contribution to bed shear stress
to interact non-linearly with the current-induced contribution, so that neither
can be obtained from the other. In the Mar Menor, a restricted micro-tidal
lagoon of comparable depth, \citet{Morote2025} found locally generated waves to
produce near-bed orbital velocities sufficient to mobilise bottom sediment over
most of the basin. Neither process has been represented in previous Stagnone
simulations. The second is the combined use of spatially distributed vegetation
roughness mapped from satellite imagery, rather than assumed uniform, and
Lagrangian drifter observations as a validation target. Together, these extensions allow the relative contributions
of three-dimensional dynamics, wave forcing, and seagrass heterogeneity to
transport skill to be quantified within a controlled process-attribution
framework.

We apply the three-dimensional unstructured-mesh hydrodynamic solver Delft3D
Flexible Mesh (D-Flow FM; \citealp{Kernkamp2011}) coupled online with the
spectral wave model SWAN \citep{Booij1999}, driven by ERA5 reanalysis
meteorological forcing and CMEMS Mediterranean physical reanalysis boundary
conditions, over a nine-day period in July 2025. The seagrass canopy is
represented as a three-dimensional drag field resolved per sigma layer, with
its parameters taken from published flume and large-scale measurements of the
same species and its spatial distribution from a supervised random-forest
classification of PlanetScope multispectral imagery acquired in August 2023,
following the mapping methodology of \citet{Maltese2025}. Model performance is assessed against in-situ water level
time series at three lagoon stations and an offshore tide gauge, and against
GPS surface drifter trajectories from a field campaign concurrent with the
simulation period, using the Liu--Weisberg Lagrangian skill score and endpoint
proximity as metrics. A complete $2\times2\times2$ factorial in wave coupling,
seagrass canopy and morphodynamic bed update attributes differences in
Lagrangian skill to individual physical mechanisms and to their interactions
(Section~\ref{sec:methods_experiments}).

The paper is organised as follows. Section~\ref{sec:study_area} describes the
study area and its hydrodynamic setting. Section~\ref{sec:data_methods}
presents the observational datasets, the modelling framework, and the seagrass
mapping and drifter validation methodology. Section~\ref{sec:results} reports
validation results for water level and Lagrangian trajectories and the
process-attribution experiment. Section~\ref{sec:discussion} discusses the
implications for wind-driven vertical structure and seagrass--flow
interaction. Section~\ref{sec:conclusions} provides conclusions.
